% Part: sets-functions-relations
% Chapter: relations
% Section: reflections
%
\documentclass[../../../include/open-logic-section]{subfiles}

\begin{document}

\olfileid{sfr}{rel}{ref}
\olsection{দার্শনিক বিবেচনা}

\olref[set]{sec}-এ সম্পর্ককে বিশেষ ধরনের সেট হিসেবে সংজ্ঞায়িত করেছি।
এখানে একটু থেমে একটি দার্শনিক প্রশ্ন করা দরকার: এমন সংজ্ঞা
আসলে কী \emph{কাজ করছে}? আমরা কিছু অধিবিদ্যাগত অভিন্নতার তথ্য
\emph{আবিষ্কার করেছি}, এমন দাবি করতে চাই কি না, তা খুবই সন্দেহজনক।
যেমন, \olref[set]{sec}-এ সংজ্ঞায়িত
$R=\Setabs{\tuple{n,m}}{n, m \in \Nat\text{ এবং } n < m}$
সেটটিই যে $\Nat$-এর উপর ক্রমসম্পর্ক, এমনটা \emph{ধরা পড়েছে}
বলা চলে কি? সন্দেহের তিনটি কারণ দিচ্ছি।

প্রথমত, \olref[set][pai]{wienerkuratowski}-তে
$\tuple{a, b} = \{\{a\}, \{a, b\}\}$ সংজ্ঞায়িত করেছি।
তার পরিবর্তে এই সংজ্ঞাটি বিবেচনা করো:
$\lVert a, b\rVert = \{\{b\}, \{a, b\}\} = \tuple{b,a}$।
$a \neq b$ হলে $\tuple{a, b} \neq \lVert a,b\rVert$।
কিন্তু ক্রমযুগলের সংজ্ঞা হিসেবে $\tuple{a,b}$-এর বদলে
$\lVert a,b\rVert$-কেও একইভাবে বেছে নিতে পারতাম। দুটি সংজ্ঞাই
সমানভাবে কাজ করত। তাই স্বাভাবিক সংখ্যার উপর ক্রমসম্পর্ক ``হওয়ার''
জন্য এখন দুটি সমান উপযুক্ত প্রার্থী আছে:
\begin{align*}
		R &= \Setabs{\tuple{n,m}}{n, m \in \Nat \text{ এবং }n < m}\\
		S &= \Setabs{\lVert n,m\rVert}{n, m \in \Nat \text{ এবং }n < m}.
\end{align*}
সদস্যভিত্তিক সমতার নীতি অনুযায়ী $R \neq S$। তাই স্পষ্টতই
তারা \emph{উভয়েই} $\Nat$-এর উপর ক্রমসম্পর্কের সঙ্গে অভিন্ন হতে পারে না।
কিন্তু $S$-এর পরিবর্তে $R$-ই (অথবা উল্টোটি) প্রকৃতপক্ষে সেই ক্রমসম্পর্ক
\emph{হয়}, এমন দাবি নিছক ইচ্ছামতো নির্বাচন হবে, আর তাই কিছুটা
বিব্রতকরও হবে। (এটি সেটতত্ত্বে পর্যবসনবাদের বিরুদ্ধে এমন এক যুক্তির
খুব সরল উদাহরণ, যা বেনাসেরাফ \citeyear{Benacerraf1965}-এ সুপরিচিত
করে তুলেছিলেন। এই যুক্তিতে আমরা কয়েক বার ফিরব।)

দ্বিতীয়ত, যদি মনে করি \emph{প্রতিটি} সম্পর্ককে কোনও সেটের সঙ্গে
অভিন্ন বলে ধরতে হবে, তবে সেট-সদস্যতার সম্পর্ক $\in$-ও একটি নির্দিষ্ট
সেট হবে। সেটিকে হতে হবে $\Setabs{\tuple{x,y}}{x \in y}$।
কিন্তু এই সেট কি আছে? রাসেলের কূটাভাসের কথা মনে রাখলে বোঝা যায়,
এমন সেটের অস্তিত্ব দাবি করা সহজ কথা নয়। বস্তুত,
\oliflabeldef{cumul:::part}{\olref[cumul][][]{part}-এ আমরা যে সেটতত্ত্ব
গড়ে তুলব, তা এই সেটের অস্তিত্ব \emph{অস্বীকার} করবে।\footnote{আগের
আলোচনা কিছুটা এগিয়ে নিয়ে কারণটি বলা যায়। স্ববিরোধের মাধ্যমে প্রমাণের
জন্য ধরি, $I = \Setabs{\tuple{x,y}}{x \in y}$ আছে। তাহলে
$\bigcup \bigcup I$ হবে সার্বিক সেট; এটি
\olref[sfr][z][sep]{thm:NoUniversalSet}-এর বিরোধী।}}{সেটতত্ত্বকে
স্বীকার্যভিত্তিক তত্ত্ব হিসেবে কঠোরভাবে গড়ে তোলা যায়, এবং সেই তত্ত্ব
এই সেটের অস্তিত্ব অস্বীকার করবে।}
তাই কিছু সম্পর্ককে সেট হিসেবে ধরা গেলেও, সেট-সদস্যতার সম্পর্ককে
বিশেষ ক্ষেত্র হিসেবে বিবেচনা করতেই হবে।

তৃতীয়ত, সম্পর্ককে সেটের সঙ্গে ``অভিন্ন'' বলে ধরার সময় বলেছি যে,
$\tuple{x,y} \in R$ বোঝাতে $Rxy$ লিখতে পারব। এতে অসুবিধা নেই,
যদি সদস্যতার সম্পর্ক ``$\in$''-কে একটি বিধেয় \emph{হিসেবে} ধরা হয়।
কিন্তু যদি মনে করি ``$\in$'' কোনও বিশেষ ধরনের সেট নির্দেশ করে,
তবে ``$\tuple{x,y} \in R$'' অভিব্যক্তিটি কেবল সেট নির্দেশকারী তিনটি
একবস্তুনির্দেশক পদ নিয়ে গঠিত: ``$\tuple{x,y}$'', ``$\in$'' এবং
``$R$''। নামের এমন তালিকা কোনও বচন প্রকাশ করতে পারে না,
যেমন পারে না এই অর্থহীন প্রতীকক্রমটি: ``কাপটি কলমদানি টেবিলটি''।
আবারও দেখা যাচ্ছে, কিছু সম্পর্ককে সেট হিসেবে ধরা গেলেও,
সেট-সদস্যতার সম্পর্ককে বিশেষ ক্ষেত্র হিসেবে ধরতে হবে। (এখানে
ফ্রেগের \emph{ঘোড়া} ধারণা-সংক্রান্ত কূটাভাসের একটি সরল রূপ এবং
রাসেলের বিরুদ্ধে ভিটগেনস্টাইনের এক বিখ্যাত আপত্তি একত্র করা হয়েছে।)

তাহলে আমাদের অবস্থান কী দাঁড়াল? সংখ্যার উপর সম্পর্কগুলি সেট,
এ কথা বলায় \emph{ভুল} কিছু নেই। শুধু বুঝতে হবে, কথাটি কোন
উদ্দেশ্যে বলা হচ্ছে। আমরা কোনও অধিবিদ্যাগত অভিন্নতার তথ্য দিচ্ছি না।
কেবল লক্ষ করছি যে, নির্দিষ্ট প্রসঙ্গে কিছু সম্পর্ককে নির্দিষ্ট সেট
\emph{হিসেবে ধরতে} পারি, এবং ধরবও।

\end{document}
